\begin{abstract}
Retrieval-augmented generation systems increasingly flag contradicting retrieved passages to the generator. We ask
what such a flag must say, and how reliable its detector must be, to counter misinformation. We vary detector and
flag wording on fixed passages for seven generators. For four, a neutral flag (``these passages may conflict'') moves
answers along the prompt's answer policy: towards abstaining when it asks for one answer, towards listing both sides
when it asks for every valid answer. The other three barely move or move the other way. A flag naming the false side
helps when right, but generators also follow wrong ones: with direction randomised, three of four generators
followed wrong flags about as often as right ones, also on new questions, and gpt-6.1-terra followed
\rfolWrongTerra\%; Claude Sonnet, the most accurate there, followed \cxRdSonnetFolW\%. Such flags pay only above a
break-even direction precision that differs by generator (exploratory). An account
frozen in advance, assuming follow rates transfer, predicted effects on unseen questions with mean
absolute error \frMae\ points, but a baseline ignoring which questions are flagged scored \frMaeMarg, and the account
overstated harm on conflict-free questions. Directional results come from ConflictBank, whose synthetic
misinformation concerns facts the generators rarely recall, yet whose true answer detectors pick from the two
candidates without passages and whose false side a bag-of-words classifier reads from style; all are scored by
string matching. We derive a rule, not tested end to end, for choosing a detector or
none, and release code and runs.
\end{abstract}
